\documentclass[10pt,a4paper]{amsart}
\usepackage[margin=19mm]{geometry}
\usepackage{amsmath,amssymb,mathtools}
\usepackage[scr=rsfs,cal=euler]{mathalfa}
\usepackage{xfrac}
\usepackage[unicode,hidelinks,bookmarks=false]{hyperref}
\newcommand{\ie}{i.e.\ }
\newcommand{\cf}{cf.\ }
\newcommand{\resp}{respectively\ }
\newcommand{\Subsection}{\subsection}
\providecommand{\nobreakdash}{\nobreak\hbox{-}}
\setlength{\emergencystretch}{2em}
\multlinegap0pt
\usepackage{amssymb}
\newtheorem{theorem}{Theorem}
\def\e{\varepsilon}
\title{Dimension reduction of fractional Sobolev seminorms\\ on thin domains}
\author{Andrea Braides, Andrea Pinamonti, Margherita Solci}
\date{Source: arXiv:2603.13968v1 (CC BY 4.0)}
\begin{document}
\maketitle

\noindent\emph{Source and licence.} Dimension reduction of fractional Sobolev seminorms\\ on thin domains. Version arXiv:2603.13968v1 (CC BY 4.0), distributed by arXiv under the Creative Commons Attribution 4.0 International licence, http://creativecommons.org/licenses/by/4.0/, as recorded in the arXiv licence metadata for this exact version and shown on the abstract page https://arxiv.org/abs/2603.13968v1. Full licence notice: \texttt{licenses/arxiv-2603.13968-CC-BY-4.0.txt}.\par\smallskip

\noindent\textit{Editorial scope.} Counted unit: the article's theorem Gammaths0 (source0.tex lines 1267--1277). The article's complete original proof of it is delivered with it (source0.tex lines 1288--1386, one \texttt{proof} environment of 99 lines, which the article places away from the statement). No mathematics is written, repaired or re-derived: the statement and the proof are the article's own lines, in the article's own notation, and no retained line is substituted or re-flowed.  Retained uncounted as the source-local context the proof needs: lines 979--982.  Nothing was omitted from any retained span, so the package carries no omission record.  No retained line contains a citation, so the wrapper carries no bibliography and no bibliography entry.  No retained line contains a figure environment, an image-inclusion command or drawing code, so no image is delivered and the package declares no asset dependency.  Editorial formatting only, in addition: an AMS \texttt{amsart} preamble that restates the article's own declarations and macros used by the retained lines, namely the environments \texttt{theorem} on separate counters, together with the standard packages those lines need.  The retained environments print in this document as Theorem 1 in order of appearance, whereas the article prints the same items with the numbering of its own full document.  The complete native source of the article as distributed in the arXiv e-print for this exact version is bundled unchanged as \texttt{sources/196364/source0.tex}.
\begin{theorem}[Non-local dimension-reduction compactness]\label{nlcth-1}
Let $s=s_\e\to s_0>1/2$ and let $u_\e$ be such that $$\sup_\e\frac{1-s}{\e^{3-2s}}\lfloor u_\e\rfloor_s^2(\Omega_\e)\le S<+\infty.$$Then there exist $u\in H^1(\omega)$ and a subsequence $\{u_{\e_j}\}_j$ such that, up to the addition of constants, 
$u_{\e_j}\buildrel{\hbox {\rm\scriptsize DR }\, \, }\over{\longrightarrow} u$ locally in $L^1$.
\end{theorem}

\begin{theorem}[Dimension-reduction Gamma-limit in the super-critical regime]\label{Gammaths0} Let $s=s_\e\to s_0\in(1/2,1)$ as $\e\to 0$. 
Then we have
$$
\Gamma\hbox{-}\lim_{\e\to 0} \frac{1}{\e^{3-2s_\e}}\lfloor u\rfloor_{s_\e}^2(\Omega_\e)= \frac1{1-s_0} K_{s_0,d}\int_\omega |\nabla'u|^2dx'
$$
for all $u\in H^1(\omega)$, where the $\Gamma$-limit is computed with respect to the dimension-reduction convergence in $L^1_{\rm loc}$, and
\begin{equation}
K_{s,d}=  \frac{1}{(3-2s)(d-1)} \int_{
 \mathbb R^{d-1}}\frac{|\xi'|^2}{(1+|\xi'|^2)^{\frac d2+s}}d\xi'.
\end{equation}
\end{theorem}

\begin{proof}[Proof of Theorem \rm\ref{Gammaths0}] We use the notation
$$
F^s_\e(u)=\frac{1}{\e^{3-2s}}\lfloor u\rfloor_{s}^2(\Omega_\e)
$$
{\em Lower bound.}
Let $u_\e\buildrel{\hbox {\rm\scriptsize DR }\, \, }\over{\longrightarrow} u\in H^1(\omega)$ in $L^1_{\rm loc}$. Up to restricting the lower bound to a slightly smaller thin film, we can suppose that the convergence of the functions $v_\e$, where $v_\e(x',x_d)= u_\e(x', \e x_d)$, to $u$ is in $L^2(\omega\times (0,\e))$. With fixed $K>0$ for each $\omega'$ compact subset of $\omega$ we can estimate
\begin{eqnarray}
\label{newest}
F^s_\e(u_\e)\ge\frac{1}{\e^{3-2s}}\int_{B_{K\e}}\int_{\omega'}
\int_0^\e\int_0^\e \frac{|u_\e(x'+\xi',y_d)-u_\e(x',x_d)|^2}{((x_d-y_d)^2+|\xi|^2)^{\frac d2+s}}dx' dx_d dy_d d\xi'
\end{eqnarray}
for $\e$ small enough.
For every $\xi'\neq 0$ we introduce the probability measures
$$
d\mu_\e^{\xi'}=\frac{1}{C_\e(\xi') ((x_d-y_d)^2+|\xi'|^2)^{\frac d2+s}}dx_d dy_d,
$$
where
\begin{equation}\label{defCexi}
C_\e(\xi')=\int_0^\e\int_0^\e\frac{1}{((x_d-y_d)^2+|\xi'|^2)^{\frac d2+s}}dx_d dy_d,
\end{equation}
and the averaged functions
$$
u^{\xi'}_\e(z)=\int_0^\e\int_0^\e u_\e(z,x_d)d\mu_\e(x_d,y_d)=\int_0^\e\int_0^\e u_\e(z,y_d)d\mu_\e(x_d,y_d).
$$
Note that we have $u^{\xi'}_\e\to u$ in $L^2(\omega)$. From Jensen's inequality we then deduce that
\begin{eqnarray*}
F^s_\e(u_\e)\ge\frac{1}{\e^{3-2s}}\int_{B_{K\e}}|\xi'|^2 C_\e(\xi') \int_{\omega'}
\frac{|u^{\xi'}_\e(x'+\xi')-u^{\xi'}_\e(x')|^2}{|\xi'|^2} dx' d\xi'.
\end{eqnarray*}
We use the notation $\overline u_\e^{\xi'}$ for the piecewise-affine interpolations on lattices aligned with the vector $\xi'$ (in $\mathbb R^{d-1}$ instead of $\mathbb R^d$), we then obtain 
\begin{eqnarray}\label{corpo}
F^s_\e(u_\e)\ge\frac{1}{\e^{3-2s}(d-1)}\int_{B_{K\e}}|\xi'|^2 C_\e(\xi') \int_{\omega''}
|\nabla'\overline u_\e^{\xi'}|^2 dx' d\xi',
\end{eqnarray}
up to restricting to a slightly smaller $\omega''$. Up to a further average, analogous to that in Step 5 of the proof of Theorem \ref{nlcth-1}, we can suppose that $\overline u_\e^{\xi'}\to u$ weakly in $H^1(\omega'')$, 
so that 
\begin{eqnarray}\label{newest-2}
\liminf_{\e\to 0}F^s_\e(u_\e)\ge\liminf_{\e\to 0}\frac{1}{\e^{3-2s}(d-1)}\int_{B_{K\e}}|\xi'|^2 C_\e(\xi') d\xi'
\int_{\omega''} |\nabla' u|^2 dx' .
\end{eqnarray}
By a change of variable and Fubini's Theorem, we can compute
\begin{eqnarray*}
\int_{B_{K\e}}|\xi'|^2 C_\e(\xi') d\xi'
&=&\int_0^\e\int_0^\e\int_{B_{K\e}}\frac{|\xi'|^2}{((x_d-y_d)^2+|\xi'|^2)^{\frac d2+s}}d\xi'dx_d dy_d
\\
&=&\e^{3-2s}\int_0^1\int_0^1|t-\tau|^{1-2s}\int_{B_{\frac{K}{|t-\tau|}}}\frac{|\xi^\prime|^2}{(1+|\xi^\prime|^2)^{\frac d2+s}}d\xi^\prime dt d\tau. 
\end{eqnarray*}
Noting that, by Lebesgue Dominated Convergence Theorem,  
\begin{eqnarray*}
&&\hskip-1cm\lim_{K\to+\infty}\int_0^1\int_0^1|t-\tau|^{1-2s}\int_{B_{\frac{K}{|t-\tau|}}}\frac{|\xi^\prime|^2}{(1+|\xi^\prime|^2)^{\frac d2+s}}d\xi^\prime dt d\tau\\
&&=\int_0^1\int_0^1|t-\tau|^{1-2s} dt d\tau\int_{\mathbb R^{d-1}}\frac{|\xi^\prime|^2}{(1+|\xi^\prime|^2)^{\frac d2+s}}d\xi^\prime \\
&&=
\frac1{(1-s)(3-2s)}\int_{\mathbb R^{d-1}}\frac{|\xi^\prime|^2}{(1+|\xi^\prime|^2)^{\frac d2+s}}d\xi^\prime, 
\end{eqnarray*}
we obtain
\begin{equation}\label{limK}
\lim_{K\to+\infty} \int_{B_{K\e}}|\xi'|^2 C_\e(\xi') d\xi'
=\frac{\e^{3-2s}}{(1-s)(3-2s)}\int_{\mathbb R^{d-1}}\frac{|\xi^\prime|^2}{(1+|\xi^\prime|^2)^{\frac d2+s}}d\xi^\prime. 
\end{equation}
Eventually, we then have the estimate 
$$
\liminf_{\e\to 0}F^s_\e(u_\e)\ge \frac{1}{1-s_0}K_{s_0,d}\int_{\omega''} |\nabla' u|^2 dx',
$$
and we can finally let $\omega''$ invade $\omega$. 
\smallskip

{\em Pointwise convergence and upper bound.}
We show a pointwise convergence result for $u\in C^2(\mathbb R^{d-1})$. 
%{\color{red} qui si usa $\omega$ limitato, ma viene la stessa cosa se assumiamo invece $u$ costante all'infinito} 

Note that, if we fix $K>0$ arbitrary, then the contribution of the integral in the set $\{|x^\prime-y^\prime|\geq \e K\}$ is negligible as $\e\to 0$ with respect to the contribution in the complementary set. Indeed,
\begin{eqnarray*}
\int_{\{|x^\prime-y^\prime|\geq \e K\}} \frac{|u(x^\prime)-u(y^\prime)|^2}{|x-y|^{d+2s}}\, dx\, dy\leq 
L^2|\omega|\int_{\{\e K\leq |\xi^\prime|\}} C_\e(\xi^\prime)|\xi^\prime|^2 d\xi^\prime, 
\end{eqnarray*}
where $L$ is a Lipschitz constant for $u$ in $\omega$, and $C_\e(\xi')$ is defined in \eqref{defCexi}. 
Since $C_\e(\xi')\leq \e^{2}|\xi^\prime|^{-d-2s}$ and $s>\frac{1}{2}$, we obtain  
\begin{eqnarray*}
\frac{1}{\e^{3-2s}}\int_{\{|x^\prime-y^\prime|\geq \e K\}} \frac{|u(x^\prime)-u(y^\prime)|^2}{|x-y|^{d+2s}}\, dx\, dy &\leq &C\frac{1}{\e^{1-2s}}
\int_{\{|x^\prime-y^\prime|\geq \e K\}} |\xi^\prime|^{2-d-2s} d\xi^\prime\\
&=&C\frac{1}{\e^{1-2s}} \sigma_{d-1}\int_{\e K}^{+\infty} t^{-2s}\, dt\\
&=&C\frac{1}{2s-1} \sigma_{d-1}K^{1-2s}, 
\end{eqnarray*} 
where $C>0$ does not depend on $K$, $\e$ and $s$. 
Since $K$ can be chosen arbitrarily large and $1-2s<0$, we obtain that we can consider the $s$-seminorm only in the set $\{|x^\prime-y^\prime|<\e K\}$. 
Then, we can estimate
\begin{eqnarray*}
{F_\e^{s}(u)}\leq \frac{1}{\e^{3-2s}}\int_{B_{K\e}}C_\e(\xi^\prime)\int_{\omega}
|u(x'+\xi')-u(x')|^2 dx^\prime d\xi^\prime+o(1)_{K\to+\infty}, 
\end{eqnarray*}
obtaining by symmetry 
\begin{equation}\label{up}
{F_\e^{s}(u)}\leq 
\frac{1}{\e^{3-2s}(d-1)}\int_{B_{K\e}}|\xi'|^2 C_\e(\xi') d\xi^\prime \int_{\omega}
|\nabla u|^2 dx^\prime +o(1)_{K\to+\infty}.
\end{equation}
Since \eqref{limK} holds, the upper bound follows
by letting $K\to+\infty$ in \eqref{up}, and recalling the lower-bound estimate we obtain the pointwise convergence of $F_\e^s(u)$. 
\end{proof}

\end{document}
